% generated by Paperino/scripts/faults.py -- do not edit by hand
\begin{table}[H]
\caption{Fault injection on the deployed speed layer.}
\label{tab:faults}
\centering
\small
\setlength{\tabcolsep}{3pt}
\begin{tabular}{lccccccc}
\toprule
 & \multicolumn{2}{c}{Windows} & & \multicolumn{2}{c}{Duplicated} & & \\
\cmidrule(lr){2-3}\cmidrule(lr){5-6}
Scenario & Prec. & Rec. & Lost & Log & Files & Starts & Recovery (s) \\
\midrule
No fault & 1.000 & 1.000 & 0 & 0 & 0 & -- & -- \\
\texttt{kill -9} of the Spark JVM & 1.000 & 1.000 & 0 & 0 & 0 & 1--2 & 18--32 \\
Kafka broker restart & 1.000 & 1.000 & 0 & 0 & 0 & 0 & 10 \\
Offsets lost (retention) & 0.965 & 0.917 & 4,460--4,474 & 0 & 0 & 1 & 5--6 \\
Duplicate block re-sent & 0.967 & 0.967 & 0 & 200 & 200 & -- & -- \\
Block behind the watermark & 0.983 & 0.983 & 0 & 0 & 0 & -- & -- \\
\bottomrule
\end{tabular}
\\[2pt] \footnotesize Each row gives the range over 3 runs. One active archive hour (55,127 events) replayed in real time into a private broker; the deployed job (two queries, Parquet sinks, checkpoints, \texttt{failOnDataLoss=false}) runs under the restart loop of the deployment. Precision and recall are over one-minute windows: a speed-layer window is correct when all four counts equal the batch ground truth. Lost and duplicated events are counted by sequence number in the archive sink, as Spark's commit log records it (log) and as the batch layer and API read it, every part file (files). Starts counts job processes the restart loop launched after the fault. Recovery is the time from the fault (for retention, from the consumer restart) to the commit of the first metrics batch planned after it that read new input, in the job process that then stayed up. Host sleep was not tested.
\end{table}
